\documentclass[10pt,twocolumn]{article}
\usepackage[utf8]{inputenc}
\usepackage[spanish,es-nodecimaldot]{babel}
\usepackage{amsmath,amssymb}
\usepackage{graphicx}
\usepackage[margin=1.65cm,columnsep=0.65cm]{geometry}
\usepackage{url}

\setlength{\parindent}{0.9em}
\setlength{\parskip}{0.15em}
\newcommand{\LCDM}{$\Lambda$CDM}
\newcommand{\Msun}{M_{\odot}}
\newcommand{\Mearth}{M_{\oplus}}
\newcommand{\Rearth}{R_{\oplus}}

\title{\textbf{Demostración empírica de capacidad explicativa multiescala}\\
\large Control cosmológico y validación cruzada en la relación masa-radio planetaria}
\author{Germán García\\\small Investigador independiente}
\date{15 de agosto de 2026}

\begin{document}
\maketitle

\begin{abstract}
Se somete un operador multiescala a dos pruebas empíricas con resultados
opuestos. En Pantheon+, una transición cosmológica de un parámetro reduce
ligeramente el residuo, pero no supera la penalización bayesiana ni produce
una ganancia predictiva material. En 897 exoplanetas con masa y radio
medidos, una ley continua de tres regímenes supera a una potencia única y a
un polinomio cuadrático tanto en BIC como en validación cruzada de cinco
pliegues. El resultado demuestra capacidad descriptiva, predictiva y
selectiva del tratamiento por escalas en el problema planetario. No
establece una ley fundamental, no demuestra unificación física y no convierte
las fronteras inferidas en constantes universales.
\end{abstract}

\section{Pregunta y criterio}

La pregunta operacional no es si una curva flexible puede ajustarse a los
datos, sino si una coordenada común de escala permite identificar cambios de
régimen que sobrevivan tres controles: penalización por complejidad,
predicción fuera de muestra y estabilidad matricial. Se usa
\begin{equation}
 s=\ln\lambda, \qquad
 O(s)=O(0)\exp\left[\int_0^s \alpha_O(u)\,du\right],
 \label{eq:general}
\end{equation}
donde $\alpha_O=d\ln O/ds$ es un exponente local. Un modelo multiescala debe
reducirse a un cierre simple cuando $\alpha_O$ es constante y permitir una
transición suave cuando cambia el régimen.

Para cada ajuste se calcula el operador linealizado y normalizado
\begin{equation}
 C_{\mathrm{eff}}=W^{1/2}\,J\,D^{-1},
 \qquad J_{ij}=\frac{\partial \mu_i}{\partial\theta_j},
\end{equation}
donde $W$ es la precisión observacional y $D$ normaliza las columnas. Su
número de condición $\kappa=\sigma_{\max}/\sigma_{\min}^{+}$ mide
identificabilidad local. Un residuo pequeño con $\kappa$ grande puede cerrar,
pero es menos robusto.

\section{Control negativo: Pantheon+}

\subsection{Datos y modelos}

Se usaron 1,590 filas de Pantheon+ con $z_{\mathrm{HD}}>0.01$ y la matriz
oficial completa de covarianza estadística y sistemática. La selección y la
teoría de distancia reproducen el módulo cosmológico público del equipo
Pantheon+ \cite{Brout2022,DataRelease}. Las observaciones se separaron por
identificador físico CID: 1,270 filas para entrenamiento y 320 para prueba.
La puntuación reservada usa la distribución gaussiana condicional, incluyendo
la covarianza entre ambas particiones.

El modelo base es
\begin{equation}
 E^2(z)=\Omega_m(1+z)^3+(1-\Omega_m).
\end{equation}
Se comparó con $w$CDM y con una transición multiescala de igual número de
parámetros:
\begin{align}
 \frac{d\ln\rho_X}{ds}&=\alpha_X(s)=A\,g(s),\\
 g(s)&=\frac{1-\tanh[(s-s_t)/\Delta_s]}{2},
\end{align}
con $s=\ln(1+z)$, $z_t=0.5$ y $\Delta_s=0.12$ fijados antes del ajuste. El
único parámetro adicional es $A$.

\begin{table}[t]
\centering
\caption{Resultados cosmológicos. Menor AICc/BIC es mejor.}
\small
\begin{tabular}{lrrrr}
\hline
Modelo & $\chi^2$ & BIC & $\chi^2_{\rm test}/N$ & $\kappa$\\
\hline
\LCDM & 1402.919 & 1417.662 & 0.862681 & 3.38\\
$w$CDM & 1402.467 & 1424.581 & 0.862115 & 16.95\\
Multiescala & 1402.424 & 1424.539 & 0.862047 & 11.99\\
\hline
\end{tabular}
\label{tab:cosmo}
\end{table}

\subsection{Resultado}

La amplitud inferida fue $A=0.267\pm0.342$ y, por tanto, compatible con cero
en el diagnóstico linealizado. Frente a \LCDM, la extensión obtiene
$\Delta\mathrm{BIC}=+6.877$ y sólo mejora $\chi^2$ reservado en 0.203. El
criterio preregistrado exigía $\Delta\mathrm{BIC}<-2$, una mejora reservada
mayor que 2 y superar a $w$CDM. No se cumple. Pantheon+ actúa así como control
negativo: el operador no obliga a declarar una transición cuando los datos no
la justifican.

\begin{figure*}[t]
\centering
\includegraphics[width=0.96\textwidth]{figura_demostracion_capacidad_multiescala.png}
\caption{Prueba cosmológica con covarianza completa. La reducción residual de
la transición multiescala es demasiado pequeña para compensar el parámetro
adicional. El pie incorporado en la figura documenta muestra, partición y
alcance.}
\label{fig:cosmo}
\end{figure*}

\section{Prueba positiva: exoplanetas}

\subsection{Muestra y verosimilitud}

Se usó una instantánea fija del NASA Exoplanet Archive del 13 de mayo de 2023
\cite{Akeson2013,NASATAP}, preservada con commit y hash. De 5,347 filas se
seleccionaron 897 planetas transitantes con masa física, no $M\sin i$, radio,
banderas de límite nulas, error relativo máximo de 50\% en masa y 20\% en
radio. Cada planeta fue prueba exactamente una vez en cinco pliegues
estratificados por masa.

Con $x=\ln(M/\Mearth)$ y $y=\ln(R/\Rearth)$ se compararon:
\begin{align}
 y_{\rm pot} &= a+b(x-\ln 10),\\
 y_{\rm cuad} &= a+b x_c+c x_c^2,\\
 y_{\rm multi} &= a+b_1x_c+(b_2-b_1)h(x-x_1)\nonumber\\
 &\quad +(b_3-b_2)h(x-x_2),
\end{align}
donde $x_c=x-\ln 10$ y
$h(u)=\Delta\ln[1+\exp(u/\Delta)]$, con $\Delta=0.16$. Las dos fronteras
$x_1,x_2$ y los tres exponentes se ajustaron, no se impusieron. La varianza de
cada residuo fue
\begin{equation}
 \sigma_i^2=\sigma_{y,i}^2+
 \left(\frac{dy}{dx}\right)^2\sigma_{x,i}^2+\sigma_{\rm int}^2,
\end{equation}
por lo que la inferencia incluye incertidumbre en ambos ejes y dispersión
intrínseca. Las leyes rotas probabilísticas ya son una clase establecida para
masa-radio planetaria \cite{ChenKipping2017,Wolfgang2016}; esta prueba evalúa
si el formalismo multiescala recupera esa estructura, no reclama haberla
descubierto por primera vez.

\begin{table}[t]
\centering
\caption{Resultados masa-radio. NLL-CV predice cada planeta una vez.}
\small
\begin{tabular}{lrrrr}
\hline
Modelo & BIC & NLL-CV/$N$ & Error med. & $\kappa$\\
\hline
Potencia & 822.412 & 0.4511 & 26.1\% & 2.49\\
Cuadrático & 455.383 & 0.2451 & 19.2\% & 6.67\\
Multiescala & 168.141 & 0.0884 & 15.5\% & 28.49\\
\hline
\end{tabular}
\label{tab:planets}
\end{table}

\subsection{Resultado}

El modelo multiescala mejora el BIC en 654.272 unidades respecto a la potencia
única y en 287.243 respecto al cuadrático. En validación cruzada reduce la NLL
por planeta de 0.451 a 0.088 y el error mediano fraccional de radio de 26.1\%
a 15.5\%. La mejora no depende de un único conjunto reservado: cada observación
se predice fuera de su ajuste.

El ajuste completo da exponentes
\begin{equation}
 (b_1,b_2,b_3)=(0.442,\;0.730,\;-0.015)
\end{equation}
y fronteras en $19.6\,\Mearth$ y $118.7\,\Mearth$. Sin embargo, la estabilidad
entre pliegues distingue ambas: la inferior es
$14.9\pm11.8\,\Mearth$, mientras que la superior es
$128.0\pm21.6\,\Mearth$. La primera no es una frontera robusta. Además,
$\kappa=28.49$ muestra mayor fragilidad paramétrica que en los modelos simples.
La ganancia predictiva es real dentro de esta muestra, pero la interpretación
fina de todos sus parámetros no tiene la misma solidez.

\begin{figure*}[t]
\centering
\includegraphics[width=0.96\textwidth]{figura_demostracion_exoplanetas_multiescala.png}
\caption{Prueba masa-radio con cinco pliegues. La ley por regímenes mejora la
predicción reservada y los criterios de información. La frontera alta es más
estable que la baja.}
\label{fig:planets}
\end{figure*}

\section{Dictamen acotado}

Los dos resultados juntos son más informativos que un único ajuste positivo:
\begin{enumerate}
\item \textbf{Capacidad descriptiva: sí.} La forma multiescala comprime la
relación masa-radio mejor que una potencia o curvatura global.
\item \textbf{Capacidad predictiva: sí, en exoplanetas.} La mejora persiste en
cinco pliegues fuera de muestra.
\item \textbf{Selectividad: sí.} El mismo criterio rechaza la transición
cosmológica examinada; el marco no mejora todo por construcción.
\item \textbf{Estabilidad total: no.} Una frontera planetaria es inestable y el
operador multiescala tiene mayor número de condición.
\item \textbf{Unificación fundamental: no evaluada.} Estas pruebas no incluyen
cuantización de gravedad, amplitudes de dispersión cuántica ni una derivación
de relatividad general.
\end{enumerate}

La afirmación sostenible es precisa: \emph{el operador por exponentes locales
y transiciones de escala tiene capacidad explicativa y predictiva demostrada
para un sistema físico con regímenes observacionales, y sabe producir un
resultado nulo en un control cosmológico}. Su novedad eventual depende de
mostrar ventajas frente a métodos de cambio de régimen ya existentes y de
extender la validación a conjuntos independientes.

\section{Reproducibilidad}

Los scripts ejecutables, datos congelados, hashes, tablas y JSON de salida se
entregan junto al informe. Los comandos principales son:
\begin{verbatim}
python demostracion_capacidad_multiescala.py
python demostracion_exoplanetas_multiescala.py
\end{verbatim}
Ambos programas fijan la selección, la partición y el criterio antes de emitir
el dictamen. Los resultados numéricos completos se conservan en
\path{demostracion_multiescala_resultados.json} y
\path{demostracion_exoplanetas_resultados.json}.

\begin{thebibliography}{9}
\bibitem{Brout2022}
D. Brout et al., ``The Pantheon+ Analysis: Cosmological Constraints'',
\emph{Astrophysical Journal} 938, 110 (2022).
\url{https://doi.org/10.3847/1538-4357/ac8e04}

\bibitem{DataRelease}
Pantheon+ and SH0ES Collaboration, public data release.
\url{https://github.com/PantheonPlusSH0ES/DataRelease}

\bibitem{Akeson2013}
R. L. Akeson et al., ``The NASA Exoplanet Archive: Data and Tools for
Exoplanet Research'', \emph{PASP} 125, 989 (2013).
\url{https://doi.org/10.1086/672273}

\bibitem{NASATAP}
NASA Exoplanet Archive, Table Access Protocol documentation.
\url{https://exoplanetarchive.ipac.caltech.edu/docs/TAP/usingTAP.html}

\bibitem{ChenKipping2017}
J. Chen and D. M. Kipping, ``Probabilistic Forecasting of the Masses and
Radii of Other Worlds'', \emph{Astrophysical Journal} 834, 17 (2017).
\url{https://doi.org/10.3847/1538-4357/834/1/17}

\bibitem{Wolfgang2016}
A. Wolfgang, L. A. Rogers and E. B. Ford, ``Probabilistic Mass-Radius
Relationship for Sub-Neptune-Sized Planets'', \emph{Astrophysical Journal}
825, 19 (2016). \url{https://doi.org/10.3847/0004-637X/825/1/19}
\end{thebibliography}

\end{document}
